\usepackage[utf8]{inputenc}
%%%%%	Standardpakker og -opsætning til 'normalt' dokument:
	\usepackage[utf8]{inputenc}
	\usepackage[T1]{fontenc}
	\usepackage[danish]{babel}
	\usepackage{amsmath, amsfonts, amssymb, amsthm}
	\usepackage[left=2cm, right=2cm, top=3cm, bottom=3cm]{geometry} 
	\usepackage{graphicx}

%sidehovede
\usepackage{fancyhdr}
	\pagestyle{fancy}
	\setlength{\headheight}{15pt}
\fancyhead[L]{Lucas, Carl-Emil, Lasse}
\fancyhead[C]{Group: xx}
\fancyhead[R]{5. Semester}

%landscape for images
\usepackage{rotating}
\usepackage{float}

%%% Standardoplysninger (til metadata og \maketitle):
	\author{\textbf{Group: xx} \\
    Lasse Arpe Kristensen (lakri18)\\
    Carl-Emil Dons Christensen (carlc23)\\
    Lucas Bergholdt Hansen (lucha23)
    }
	\title{ \vspace*{-2em} DM585 Distributed Web Application Project}
  
	\date{21-12-2025}
    
%%%%%	Normale basisjusteringer
	\setlength{\parindent}{0.0em}
	\setlength{\parskip}{1em}
	\linespread{1.1}
%%%%%	Comments
\usepackage{comment}

%%%%% Wrap figures
\usepackage{wrapfig}
% \usepackage[export]{adjustbox} %%% aligning images to right. https://www.overleaf.com/learn/latex/Positioning_images_and_tables

%%%%%	Ekstra pakker og indstillinger, der gør dokumentet lidt lækrere:
\usepackage{times}
%\usepackage{lastpage}
\usepackage{titlesec}
	\titlespacing{\paragraph}{0pt}{2pt}{6pt}
	\titlespacing{\section}{0pt}{8pt plus 2pt minus 2pt}{0pt}
	\titlespacing{\subsection}{0pt}{6pt plus 1pt minus 1pt}{0pt}

\begin{comment}
Følgende er kommenteret ud da det giver bogstaver til ordered lists. -CE

\usepackage{enumitem} 
	\setlist{itemsep=0.8pt}
	\setlist{topsep=0.0pt}
	\setlist[description]{leftmargin=4em,labelindent=2em}
	\setlist[enumerate,1]{label={(\alph*)}}
\end{comment}

%% Programming snippets
\usepackage{listings}


%general packages for showing codes from programmes
\usepackage{tcolorbox}
\newtcbox{\code}{on line,             % Extra option to make text stay on line instead of moving all text before and after the \code{} text to new lines.
colupper = black!95!white,            % Font color
colback  = white!85!black,            % Background color
boxrule  = 0pt,                       % Border length
left     = 1pt, 
right    = 1pt, 
top      = -1pt, 
bottom   = -1pt,                      % Extra space to the horisontal and vertical directions
arc      = 4pt}                       % The arc of the corner. You could also call it the roundness factor to the corners.

% Nu lidt kode miljø indstillinger.
% Først skal vi bruge nogle farver, som ikke kun behøves at bruges a listlings biblioteket.
\usepackage{xcolor}
\definecolor{comment_color}{rgb}{0,0.6,0}
\definecolor{number_color}{rgb}{0.2, 0.2, 0.8}
\definecolor{string_color}{rgb}{0.58,0.0,0.82}
\definecolor{code_background_color}{rgb}{0.95, 0.93, 0.90}
\definecolor{keyword_color}{rgb}{1.0,0.4,0.0}

% Nu kan vi indstille biblioteket Listings

\lstdefinestyle{mystyle}{
  backgroundcolor = \color{code_background_color},
  commentstyle = \color{comment_color},
  keywordstyle = \color{keyword_color},
  numberstyle = \sffamily\tiny\color{number_color},
  stringstyle = \color{string_color},
  basicstyle = \sffamily\footnotesize, %font families: sffamily, ttfamily
  breakatwhitespace = false,
  breaklines = true,
  prebreak=\raisebox{0ex}[0ex][0ex]{\ensuremath{\hookleftarrow}},
  captionpos = b,
  keepspaces = true,
  numbers = left,
  numbersep = 10pt,
  showspaces = false,
  showstringspaces = false,
  showtabs = true,
  tabsize = 2,
  frame = leftline
}

\lstset{style=mystyle}


\definecolor{eclipseStrings}{RGB}{0,128,0}
\definecolor{eclipseKeywords}{RGB}{0,0,0}
\colorlet{numb}{red!80!black}

\lstdefinelanguage{json}{
    basicstyle=\normalfont\ttfamily,
    commentstyle=\color{eclipseStrings}, % style of comment
    stringstyle=\color{eclipseKeywords}, % style of strings
    numbers=left,
    numberstyle=\scriptsize,
    stepnumber=1,
    numbersep=8pt,
    showstringspaces=false,
    breaklines=true,
    frame=lines,
    backgroundcolor=\color{code_background_color}, %only if you like
    string=[s]{"}{"},
    comment=[l]{:\ "},
    morecomment=[l]{:"},
    literate=
        *{0}{{{\color{numb}0}}}{1}
         {1}{{{\color{numb}1}}}{1}
         {2}{{{\color{numb}2}}}{1}
         {3}{{{\color{numb}3}}}{1}
         {4}{{{\color{numb}4}}}{1}
         {5}{{{\color{numb}5}}}{1}
         {6}{{{\color{numb}6}}}{1}
         {7}{{{\color{numb}7}}}{1}
         {8}{{{\color{numb}8}}}{1}
         {9}{{{\color{numb}9}}}{1}
}


%YAML highligthing fra: https://tex.stackexchange.com/questions/152829/how-can-i-highlight-yaml-code-in-a-pretty-way-with-listings
\newcommand\YAMLcolonstyle{\color{red}\mdseries}
\newcommand\YAMLkeystyle{\color{black}\bfseries}
\newcommand\YAMLvaluestyle{\color{blue}\mdseries}

\makeatletter

% here is a macro expanding to the name of the language
% (handy if you decide to change it further down the road)
\newcommand\language@yaml{yaml}

\expandafter\expandafter\expandafter\lstdefinelanguage
\expandafter{\language@yaml}
{
  keywords={true,false,null,y,n},
  keywordstyle=\color{darkgray}\bfseries,
  basicstyle=\YAMLkeystyle,                                 % assuming a key comes first
  sensitive=false,
  comment=[l]{\#},
  morecomment=[s]{/*}{*/},
  commentstyle=\color{purple}\ttfamily,
  stringstyle=\YAMLvaluestyle\ttfamily,
  moredelim=[l][\color{orange}]{\&},
  moredelim=[l][\color{magenta}]{*},
  moredelim=**[il][\YAMLcolonstyle{:}\YAMLvaluestyle]{:},   % switch to value style at :
  morestring=[b]',
  morestring=[b]",
  literate =    {---}{{\ProcessThreeDashes}}3
                {>}{{\textcolor{red}\textgreater}}1     
                {|}{{\textcolor{red}\textbar}}1 
                {\ -\ }{{\mdseries\ -\ }}3,
}

% switch to key style at EOL
\lst@AddToHook{EveryLine}{\ifx\lst@language\language@yaml\YAMLkeystyle\fi}
\makeatother

\newcommand\ProcessThreeDashes{\llap{\color{cyan}\mdseries-{-}-}}


\usepackage{float}
\usepackage{calc}
\usepackage{gensymb}

\newcommand{\pmat}[1]{\ensuremath{\begin{pmatrix}#1\end{pmatrix}}}
\usepackage{subcaption}

%%%%% captions
\usepackage{caption}
\captionsetup{
    format = plain,
    font = footnotesize,
    labelfont = sc
}
%

\usepackage{biblatex}
\addbibresource{referencer.bib} 

%\bibliographystyle{plain} 
\bibliography{referencer} % Entries are in the referencer.bib file

\usepackage{hyperref}
\usepackage{enumitem}

\usepackage{booktabs}
\usepackage{pgfplots}
\pgfplotsset{width=6.8cm,compat=1.9}
\usepackage{comment}

%%%%%% Til referencer
%\usepackage[backend=biber]{biblatex}
%\usepackage{cite}
%\addbibresource{referencer.bib} 


